% !TeX root = closed-systems-from-comparison-completeness.tex
% ============================================================
\chapter*{Standing Terminology}
\phantomsection
\addcontentsline{toc}{chapter}{Standing Terminology}
\label{ch:standing-terminology}
% ============================================================
This section fixes recurring terminology at the level of use.  The
formal definitions and proofs remain in the chapters cited below.

\begin{description}[leftmargin=3.6cm,style=nextline]
\item[Admissible]
A structure is admissible when it is internally recoverable from the
comparison system, descends through the relevant quotient, or is
explicitly introduced as scoped realization or boundary data.

\item[Boundary data]
Data deliberately not absorbed into the intrinsic theorem stack.  Boundary
data enter a realized reading or a quantitative completion as named
conditions, not as hidden primitives.  The remaining global boundary data
are recorded in \cref{app:conditional-completion}.

\item[Closed world]
A system whose observers, comparison protocols, reports, and realized
readings are internal to the object being modeled.  No external frame is
available as a primitive source of distinctions.

\item[Closure]
The requirement that admissible distinctions and reports be recoverable
within the internal comparison, quotient, and transport structure.  In
the foundational comparison setting, closedness is expressed by
profile-pair completeness, or rectangular completeness.

\item[Coherent report]
A report that factors through the relevant quotient or transport-invariant
structure.  A report that distinguishes only representatives of the same
physical quotient does not define physical content.

\item[Comparison completeness]
The profile-pair completeness condition governing closed comparison
worlds.  Under intrinsicality and local distinguishability it is derived
rather than separately postulated; see
\cref{thm:rectangular-completeness-intrinsicality}.

\item[Comparison world]
The primitive pair \((U,\mathcal C)\), where \(U\) is a nonempty state set
and \(\mathcal C\) is a family of intrinsic binary comparison predicates.
See \cref{def:intro-comparison-world}.

\item[Diagonal redundancy]
The redundancy generated by simultaneous action of the intrinsic symmetry
group on representative-level data.  Physical content is read after
descent through this redundancy.

\item[External primitive]
A structure used as a background source of distinctions without being
reconstructed from the internal comparison stack.  Examples include a
primitive manifold, external observer frame, or unearned state-space
linearity.  Such structures may appear later as realized data only when
their status is explicitly marked.

\item[Ghost sector]
A formal smooth or transport-readable degree of freedom that is invisible
to the intrinsic comparison tower.  The interface analysis isolates the
conditions under which these sectors are absent.

\item[Intrinsic]
Recovered from the comparison data, the quotient semantics, or the
transport obstruction hierarchy.  Intrinsic does not mean
coordinate-free presentation of an already assumed background; it means
available from within the closed stack.

\item[Local distinguishability]
The condition that genuine comparison separations are witnessed at a
finite coordinate or finite comparison stage.  It is the local witness
condition used with intrinsicality in the derivation of rectangular
completeness.

\item[Morphism locus]
The enrichment locus carried by transport, paths, loops, or holonomy.
It is one of the two possible enrichment loci in the universal
classification.

\item[Object locus]
The enrichment locus carried by representative choices over quotient
objects.  It is the other possible enrichment locus in the universal
classification.

\item[Physical quotient]
The quotient of representative data by intrinsic redundancy.  In the
basic two-factor setting it is \(\Phys=(X_A\times X_B)/G\).

\item[Realization]
A reading of intrinsic structure into smooth, Hilbert, phase,
observer-side, macroscopic, or quantitative form.  Realization statements
are theorem-level statements only relative to the realization data named
in their hypotheses.

\item[Rectangular completeness]
The closedness condition asserting profile-pair completeness for the
comparison structure.  In the present monograph it is treated as a
derived theorem under intrinsicality and local distinguishability in the
foundational comparison setting.

\item[Regular sector]
The maximal readable sector selected in the state-side reconstruction
extension under scale coherence.  It is separated from the raw limit
observable algebra, which is larger than any realized smooth reading.

\item[Semantic descent]
The passage from representative-level distinctions to quotient-level
content.  A distinction has semantic status only if it survives the
relevant descent.

\item[Transport obstruction]
The residue that remains when quotient semantics do not determine
transport around paths or loops.  The first stable visible layer of this
residue is the quadratic carrier \(F^2/F^3\).

\item[Two-locus theorem]
The classification stating that admissible enrichment beyond quotient
semantics occurs through object-level representative choice, morphism-level
transport, or both, with no third locus.
\end{description}
